\section{第一阶段：高精地图、激光雷达和“有轨电车”思维}

上一章建立路线图，本章先回到 2014--2015 年的起点。郎咸朋说，当时很多团队把自动驾驶当作“有轨电车”研发：先用高精地图在真实道路上铺一条虚拟轨道，再用大量激光雷达探测周围动态物体，最后用 if-else 和规则系统控制车辆沿着轨道行驶。今天看起来可笑，但在当时是很自然的工程选择，因为它把开放世界问题压缩成相对确定的问题。

\begin{figure}[H]
\centering
\includegraphics[width=0.94\textwidth]{hdmap-lidar-vs-vision.png}
\caption{高精地图+激光雷达 vs 视觉路线：一条依赖外部确定性，一条依赖数据和模型泛化。自制概念图，依据 00:01:32--00:04:30 与 00:25:06--00:28:46 对谈内容整理。}
\end{figure}

\begin{knowledgebox}{读图：两条路线对不确定性的处理不同}
左侧高精地图+激光雷达路线，把不确定性外包给地图、测距和规则；右侧视觉+模型路线，把不确定性交给数据和模型学习。前者短期可控，长期维护重；后者短期难训练，长期更有规模化可能。
\end{knowledgebox}

\subsection{高精地图为什么像虚拟轨道}

本节解释“有轨电车”比喻。高精地图把车道线、道路边界、红绿灯位置、路口拓扑等信息预先写进地图，相当于给车一条可预期的轨道。车辆实时行驶时，不必完全从传感器中理解世界，而是把当前位置对齐到地图，然后按规则行驶。这个方法在高速等封闭、变化慢的道路上较容易成立，但在城市普通道路上维护成本极高。

\begin{warningbox}{高精地图的根本问题是更新和覆盖}
高速公路里程有限、变化较少，适合使用高精地图；普通城市道路数量巨大、施工频繁、临时变化多，地图难以每天更新。自动驾驶如果依赖一张永远正确的地图，就会被现实世界的变化击穿。
\end{warningbox}

\subsection{激光雷达堆料为什么无法量产}

本节看早期硬件成本。访谈中提到，早期一个 64 线激光雷达可能要五六十万人民币，百度或 Cruise 的工程车一辆车装七八个激光雷达，传感器成本远高于车本身。这样的系统适合研究演示和 Robotaxi 小规模试验，但不适合几十万、几百万辆量产车。

\begin{figure}[H]
\centering
\includegraphics[width=0.86\textwidth]{sensor-cost-stack.png}
\caption{自动驾驶成本栈：早期方案被传感器、计算平台和工程车成本拖住。自制概念图，依据 00:12:07--00:13:09 对谈内容整理。}
\end{figure}

\begin{knowledgebox}{读图：研究车和量产车不是同一种约束}
图中从多颗 LiDAR、传感器套件、计算平台、工程车到量产约束。研究车可以用昂贵传感器换确定性；量产车必须考虑成本、外观、功耗、供应链、维护和用户价格。
\end{knowledgebox}

\subsection{术语消化：ODD、高精地图和激光雷达}

本节补基础概念。自动驾驶讨论里，很多术语看似常见，但如果不解释，后面 L2/L3/L4 和端到端会很难读。

\begin{knowledgebox}{术语消化：早期自动驾驶路线}
\noindent\begin{tabularx}{\textwidth}{p{0.20\textwidth}p{0.36\textwidth}X}
\toprule
术语 & 含义 & 本期中的作用 \\
\midrule
ODD & Operational Design Domain，系统被设计允许运行的区域和条件 & 早期自动驾驶试图用 ODD 穷举场景边界。 \\
高精地图 & 厘米级道路、车道、交通设施和拓扑信息 & 提供虚拟轨道，但覆盖和更新困难。 \\
LiDAR / 激光雷达 & 主动发射激光并测距的传感器 & 几何距离直接，但成本和信息密度受限。 \\
规则系统 & 工程师手写 if-else 和场景逻辑 & 可解释，但长尾场景无法穷举。 \\
\bottomrule
\end{tabularx}
\end{knowledgebox}

\subsection{本章小结}

早期自动驾驶选择高精地图和激光雷达，是因为它能把开放道路暂时变成确定系统。但这条路在城市覆盖、地图更新、成本和长尾场景上遇到天花板。下一章的 Tesla 路线，正是对这套确定性工程路线的升维替代。

